\section{Trees (Árvores Binárias)}
\label{sec:const-trees}

As construções na categoria de árvores binárias visam criar novas estruturas que respeitem a hierarquia e a lateralidade (esquerda/direita) dos nós originais. O foco principal é garantir que a propriedade de ramificação binária seja preservada durante as operações categoriais.

\subsection{Fundamentação Teórica}

\begin{itemize}
    \item \textbf{Produtos:} 
    O produto $T_1 \times T_2$ de duas árvores binárias pode ser construído como a árvore cujos nós são os pares $(u, v)$ com $u \in T_1$ e $v \in T_2$, com raiz $(\text{root}_1, \text{root}_2)$, e onde $(u,v)$ tem filho esquerdo/direito se e só se $u$ e $v$ têm filho esquerdo/direito em simultâneo.
    
    Esta construção é válida mas pode colapsar rapidamente para a árvore vazia se as duas árvores têm profundidades diferentes --- o produto nesta categoria é \emph{não trivial} e pode ser a árvore com um só nó (se as árvores têm alturas diferentes). Na prática, o produto existe mas pode ser menor do que qualquer dos fatores.

    \item \textbf{Coprodutos:} O coproduto $T_1 \sqcup T_2$ é a união disjunta das duas árvores. Para manter a estrutura de uma única árvore binária na implementação, introduz-se uma nova raiz comum cujos filhos esquerdo e direito são as raízes de $T_1$ e $T_2$, respetivamente.

    \item \textbf{Pullbacks:} O pullback de $T_1 \xrightarrow{f} T_3 \xleftarrow{g} T_2$ é a subárvore do produto $T_1 \times T_2$ consistindo nos pares $(u, v)$ tais que $f(u) = g(v)$. A estrutura de árvore (parentesco e lateralidade) é herdada diretamente do produto $T_1 \times T_2$.

    \item \textbf{Pushouts:} O pushout de $T_1 \xleftarrow{f} T_0 \xrightarrow{g} T_2$ é o coproduto $T_1 \sqcup T_2$ quocientado pela relação que identifica $f(u)$ com $g(u)$ para cada nó $u \in T_0$. Em árvores binárias, isto corresponde a ``colar'' as duas árvores ao longo da imagem comum.

    \item \textbf{Equalizadores:} O equalizador de dois morfismos $f, g: T_1 \rightrightarrows T_2$ é a subárvore de $T_1$ induzida pelo conjunto de nós $E = \{u \in T_1 : f(u) = g(u)\}$, desde que $E$ forme uma subárvore válida (deve ser fechada para progenitores e conter a raiz).

    \item \textbf{Coequalizadores:} O coequalizador de $f, g: T_1 \rightrightarrows T_2$ é o quociente de $T_2$ pela menor relação de equivalência que identifica $f(u) \sim g(u)$ e que preserva a estrutura de árvore binária (se $u \sim v$, então os seus filhos respetivos devem ser identificados).
\end{itemize}

\subsection{Implementação Computacional}

A implementação em MATLAB utiliza funções recursivas e manipulação de índices para respeitar as condições teóricas descritas.

\begin{lstlisting}[language=Matlab, caption={Produto de duas árvores binárias (representação recursiva)}]
function T = tree_product(T1, T2)
    % Constroi o produto T1 x T2 recursivamente
    % Retorna [] se um dos nos nao tem correspondente no outro
    if isempty(T1) || isempty(T2)
        T = []; return;
    end
    T.value = [T1.value, T2.value];  % par de valores
    T.left  = tree_product(T1.left, T2.left);
    T.right = tree_product(T1.right, T2.right);
end
\end{lstlisting}

\begin{lstlisting}[language=Matlab, caption={Coproduto em Trees (Nova Raiz)}]
function T = tree_coproduct(T1, T2)
    % O coproduto e formado criando uma nova raiz artificial
    T.value = []; % Raiz sem conteudo informativo direto
    T.left  = T1; % T1 torna-se o ramo esquerdo
    T.right = T2; % T2 torna-se o ramo direito
end
\end{lstlisting}

\begin{lstlisting}[language=Matlab, caption={Equalizador e Pullback em Trees}]
function [E, inc] = tree_equalizer(T1, f, g)
    % E: subconjunto dos nos de T1 onde f e g coincidem
    mask = (f == g);
    E    = find(mask);    % indices dos nos do equalizador
    inc  = E;             % inclusao em T1
    % Nota: verificar que E e fechado para progenitores
end

function [PB, pi1, pi2] = tree_pullback(T1, T2, f, g)
    % PB e a subarvore de T1 x T2 onde f(u) == g(v)
    nodes = [];
    for i = 1:T1.n
        for j = 1:T2.n
            if f(i) == g(j)
                nodes(end+1, :) = [i, j];
            end
        end
    end
    PB.nodes = nodes;
    pi1 = nodes(:,1)'; % Projecao para T1
    pi2 = nodes(:,2)'; % Projecao para T2
    % A estrutura herda-se de T1 x T2
end
\end{lstlisting}